% !TeX root = ../教材题库_全部条目_转写工作稿.tex
% Source: XJ-B1-C03-TH-P065-Y1561
% 第3章 / 3.1 函数 / 教材 p.65、66 / PDF 第72、73页
% OCR未校对：公式与文字均可能有误；未完成数学TeX转写；禁止直接发给学生
% 原题及相邻原解截图（核对用）：../题目原图/XJ-B1-C03-TH-P065-Y1561-01.webp, ../题目原图/XJ-B1-C03-TH-P065-Y1561-02.webp
% 本文件只在转写完成后作为题干使用；OCR 条目尚不是可靠题干。
\begin{bankitem}{p.65 思考（含紧邻说明）}{第3章 / 3.1 函数 / 教材 p.65、66 / PDF 第72、73页}{XJ-B1-C03-TH-P065-Y1561}{OCR未校对：公式与文字均可能有误；未完成数学TeX转写；禁止直接发给学生}
\begin{ocrdraft}
[教材p.65 / PDF第72页]
\par
？思考我们有了集合的概念，就可以将上述例子中两个变量的取值范围看成集
\par
合，将函数看成两个集合之间的对应关系．那么，上述例子中的函数有哪些共同特征？
\par
[教材p.66 / PDF第73页]
\par
上述例子中的函数有如下共同特征：
\par
（1)都包含两个非空数集（可用A，B来表示）；
\par
（2）都有一个确定的对应关系；
\par
（3）对于数集，A中的任何一个数，按照确定的对应关系，在数集B中都有唯一
\par
的数和它对应．
\par
一般地，设A，B是两个非空的实数集，如果按照某种确定的对应关系丿、，对
\par
于集合A中的任何一个数，在集合B中都有唯一的数3，和它对应，那么称这样的
\par
对应丆：为定义于A取值于B的函数，也记作
\par
3，一不）0\ensuremath{{}^\circ\mathrm C}A，3\ensuremath{{}^\circ\mathrm C}B）．
\par
其中，x叫作自变量，工的取值范围A叫作函数的定义域；与xeA对应的数3，叫作
\par
函数值，记作不工），所有函数值组成的集合\{不IxeA\}叫作函数的值域．值域是
\par
集合的子集．
\par
由函数的定义可知，函数的值域由函数的定义域和对应关系完全确定，因此确
\par
定一个函数主要取决于两个要素：定义域和对应关系．对值域的了解有助于加深对
\par
函数的认识，所以值域也是函数的要素之一
\par
要注意的是，若只给出函数的关系式3-丆0）而没有指明其定义域，那么函数
\par
的定义域就是使关系式有意义的所有x值组成的集合．
\end{ocrdraft}
\end{bankitem}
